\session{3}{10/10}{2限}{1}{生成AIのしくみと限界}

\goals{
  \item 大規模言語モデル（LLM）が文章を生成する基本的なしくみを、経営者・管理者として必要な範囲で説明できる。
  \item 生成AIの得意なこと・苦手なこと（ハルシネーション、最新情報、計算）を理解し、対策を挙げられる。
  \item プロンプトの4要素を使って、目的に合った指示を書ける。
}

\guidesec{解説}

\topic{大規模言語モデルのしくみ}
ChatGPT や Copilot の中核にあるのは、大規模言語モデル（Large Language Model、LLM）です。LLM は、
インターネット上の文章や書籍など膨大なテキストから、「ある語の並びの次にどの語が来やすいか」を学習したモデルです。
文章を生成するときは、入力（プロンプト）に続く語を、確率の高いものから選んで1つずつ付け足していきます。
この「語」の単位をトークンと呼びます。英語では単語や単語の一部、日本語では1〜2文字程度が1トークンになることが多いです。

LLM が作られる手順を簡単に示します。
\begin{enumerate}[label=\arabic*.,leftmargin=2em]
  \item \textbf{事前学習}　膨大なテキストで次のトークンを予測する訓練を行い、言語のパターンと世界に関する知識を獲得する。
  \item \textbf{指示調整（ファインチューニング）}　「質問に答える」「指示に従う」形式の例で追加学習する。
  \item \textbf{人のフィードバックによる調整}　人が好ましいと評価した出力に近づくように調整する（RLHF など）。
\end{enumerate}
重要なのは、LLM は「正しい事実のデータベースを検索している」のではなく「もっともらしい語の並びを生成している」という点です。
正しい知識が学習データに多く含まれていれば正しい出力になりやすい、というだけで、正しさが保証されているわけではありません。

\topic{得意なことと苦手なこと}
\par\noindent\begin{gtabh}{colspec={X[l] X[l]}}
  得意なこと & 苦手なこと \\
  要約、言い換え、翻訳、文体の変換 & 事実の正確さの保証（ハルシネーション） \\
  構成案やアイデアを多数出す & 学習データの締め切り以降の最新情報 \\
  定型的な文章・コードの下書き & 厳密な計算、長い数値の処理（桁の間違い） \\
  複数の視点からの検討（役割を変えて答える） & 出典を正確に示すこと（存在しない文献を作る） \\
  自然な対話での質問の掘り下げ & 自分の出力の正誤を判断すること \\
\end{gtabh}

\textbf{ハルシネーション}は、事実に基づかない内容をもっともらしく、自信のある口調で生成する現象です。
次トークン予測というしくみから必然的に生じるもので、完全にはなくせません。実在する資料名に存在しない数値を添えるなど、
半分正しい形で出ることが多いので、固有名詞・数値・文献は必ず一次資料で確かめます。

\textbf{最新情報}については、学習データに締め切り（知識カットオフ）があり、それ以降の出来事を知りません。
Web 検索や社内文書の検索結果をプロンプトに添えて答えさせる RAG（検索拡張生成）が、この弱点を補う一般的な方法です。
Copilot のように検索と連携したツールでは、出典のリンクを確認する習慣をつけてください。

\textbf{計算}については、LLM は計算機ではないため、桁の多い計算や複雑な集計で誤ることがあります。
最近のツールは計算をプログラム（Python など）に任せる機能を持っており、「計算はコードで行ってください」と指示すると精度が上がります。

\topic{経営者・管理者として知っておくべき範囲}
LLM の内部の数式を理解する必要はありません。しかし、次の点は判断に直結するので理解しておく必要があります。
\begin{itemize}
  \item 出力の正しさは保証されない。検証の工程を業務に組み込む必要がある。
  \item 入力した内容は外部のサーバーに送られる。機密情報の扱いは方針が必要（第8回・第13回）。
  \item 学習データの偏りは出力の偏りになる（第13回）。
  \item 出力の品質はプロンプトと、与える文脈（資料・データ）の品質に大きく依存する。
  \item 同じプロンプトでも出力は毎回変わる（確率的な生成）。再現性が必要な業務では設定や検証が必要。
\end{itemize}

\topic{プロンプトの基本}
プロンプトは AI に与える指示や入力文です。本講義では次の4要素を意識します（付録も参照）。
\begin{itemize}
  \item \textbf{役割}　「あなたは経営学の大学教員です」「あなたは厳しい査読者です」
  \item \textbf{目的と背景}　「経営情報の授業で、企業へのAI活用提案を作っています」
  \item \textbf{条件}　「日本国内の事例に限る」「文献名は挙げないでください」
  \item \textbf{出力の形式}　「300字以内」「10案を表で」「箇条書きで」
\end{itemize}
さらに、一度で終わらせず「なぜ」「他には」「反対の立場だと」「この案の弱点は」と対話を続けることで出力の質が上がります。
自分の考えを先に書いてから批評させると、AIの案に引きずられにくくなります。

\topic{キーワード}
\par\noindent\begin{keywords}{}
  大規模言語モデル（LLM） & 膨大なテキストから次のトークンの確率を学習し、文章を生成するモデル。 \\
  トークン & LLM が扱う文章の最小単位。単語や部分語。文脈窓や料金の単位。 \\
  ハルシネーション & 事実に基づかない内容をもっともらしく生成する現象。 \\
  知識カットオフ & 学習データの締め切り時点。それ以降の情報はモデルが持たない。 \\
  RAG（検索拡張生成） & 検索した文書をプロンプトに添えて回答させ、最新情報や社内情報を補う方法。 \\
  プロンプト & AIに与える指示・入力文。役割・目的と背景・条件・出力の形式の4要素。 \\
\end{keywords}

\guidesec{演習課題}

\exercise{1}{AI演習：同じ問いを異なる指示で投げ、出力の違いを観察する}
\instr{次の3つのプロンプトを順に Copilot に入力し、出力を並べて比べてください。何が変わったか（長さ、具体性、構成、根拠の示し方）を表にまとめ、「どの要素がどの変化を生んだか」を考察してください。}
\begin{promptbox}[A：要素なし]
中小企業が生成AIを導入するメリットを教えてください。
\end{promptbox}
\begin{promptbox}[B：役割と形式を追加]
あなたは中小企業診断士です。従業員30人の製造業が生成AIを導入するメリットを、3つの箇条書きで、各60字以内で説明してください。
\end{promptbox}
\begin{promptbox}[C：目的・背景・条件を追加]
あなたは中小企業診断士です。私は従業員30人の金属加工業の2代目経営者で、来月の経営会議で生成AIの導入を提案します。反対派の役員を説得するために、導入のメリットを3つ、各メリットには具体的な業務の例と、想定される反論への答えを添えてください。導入費用の具体的な金額は挙げないでください。表の形式で出力してください。
\end{promptbox}

\exercise{2}{ハルシネーションを検出する}
\instr{自分がよく知っている話題（出身地、アルバイト先の業界、趣味など）について、次のプロンプトで説明させ、出力の中の誤りを探してください。誤りは「完全な誤り」「半分正しい」「確かめられない」に分けて記録します。}
\begin{promptbox}[誤りを誘いやすい問い]
（話題）について、歴史的な経緯、現状を示す統計の数値、参考文献を3つ含めて400字で説明してください。
\end{promptbox}

\exercise{3}{プロンプトを改善する}
\instr{次の「悪いプロンプト」を、4要素を使って書き直してください。何を補ったかを要素ごとに説明します。}
\begin{promptbox}[悪いプロンプト]
レポートを書いて。
\end{promptbox}

\guidesec{模範解答}

\begin{answer}{演習1}
\par\noindent\begin{gtabh}{colspec={Q[l,wd=16mm] X[l] X[l] X[l]}}
  観点 & A & B & C \\
  長さ & 長い（500字以上）。一般論が並ぶ。 & 短い。制限に従う。 & 中程度。表の各セルが具体的。 \\
  具体性 & 「業務効率化」「コスト削減」など抽象的。 & 製造業向けの例（図面の整理、見積書の下書き）が入る。 & 金属加工の業務（見積、納期回答、日報）に即した例。 \\
  構成 & 見出しと箇条書きが混在。 & 3点に整理。 & メリット・業務例・反論への答えの3列。 \\
  根拠 & なし。 & なし。 & 反論への答えで根拠が示される。費用の金額は（指示通り）出ない。 \\
\end{gtabh}
考察：役割は専門性と用語の水準を、出力の形式は長さと構成を、目的と背景は具体性と「誰を説得するか」の視点を決めます。
条件（金額を挙げない）は、確かめにくい数値を最初から出させないことで、検証の手間を減らします。
最も出力の質を変えたのは「目的と背景」で、これは自分で考えないと書けない要素です。
\end{answer}

\begin{answer}{演習2}
例：出身地（人口約8万人の地方都市）について
\begin{itemize}
  \item 完全な誤り：「市の中心部に県立大学のキャンパスがある」（実際には隣の市）。
  \item 半分正しい：「人口は2020年の国勢調査で約9万5千人」（国勢調査の年は正しいが、数値は実際より多い）。
  \item 確かめられない：参考文献「○○市史編さん委員会『○○市史 第3巻』（1998）」（市史は実在するが、巻数と年が確認できない）。
\end{itemize}
気づき：地名や施設名など「存在してもおかしくない」固有名詞を、もっともらしく組み合わせてくる。
数値は桁が合っている分、かえって見抜きにくい。自分が知らない話題では、これらの誤りに気づけない。
\end{answer}

\begin{answer}{演習3}
改善例：
\begin{promptbox}[改善後]
あなたは経営学部の大学教員です（役割）。私は大学2年生で、「経営情報」の授業の課題として、コンビニエンスストアにおける情報システムの役割について1,200字のレポートを書きます（目的と背景）。まずレポートの構成案を、序論・本論・結論の見出しと各節の要点の箇条書きで提案してください。文献名は挙げないでください。本文は書かないでください（条件・出力の形式）。
\end{promptbox}
補った要素：役割（水準の設定）、目的と背景（授業の課題、文字数）、条件（文献名を挙げない、本文を書かない）、
出力の形式（構成案の箇条書き）。「本文を書かない」という条件は、AI Policy の「AIの出力をそのまま提出物にしない」を守るためのものです。
\end{answer}

\guidesec{出席確認クイズの解説}
講義サイトの出席確認ページ（第3回）の5問です。
% 出席確認クイズ 第03回　生成AIのしくみと限界
%   attendance/attendance03.html から生成（解説は本ガイド用に加筆）。

\quizq{1}{大規模言語モデル(LLM)が文章を生成する基本的なしくみはどれでしょうか?}%
  {辞書の意味を順番に並べる}%
  {人手で記述した文法規則をすべて組み合わせる}%
  {人が書いた文章をデータベースからそのまま検索して表示する}%
  {\correct{大量のテキストから学習し、次に来る語(トークン)の確率を予測して生成する}}%
  {正解は (d)。大規模言語モデル（LLM）は膨大なテキストから学習したパターンをもとに、次に来るトークンの確率を予測しながら文章を1語ずつ生成します。(c) のようにデータベースを検索して表示しているわけではないため、出典が存在しない内容でも流暢に生成してしまいます。これがハルシネーションの根本的な原因です。(a)(b) は古典的な手法の説明です。}

\quizq{2}{「ハルシネーション」の説明として適切なものはどれでしょうか?}%
  {AIの動作が遅くなること}%
  {\correct{AIが事実に基づかない内容をもっともらしく生成すること}}%
  {AIが同じ回答を繰り返すこと}%
  {AIが画像を認識できないこと}%
  {正解は (b)。ハルシネーションは、事実に基づかない内容をもっともらしく、自信のある口調で生成してしまう現象です。(a)(c)(d) は別の問題です。実在する資料名に存在しない数値を添えるなど「半分正しい」形で現れることが多く、固有名詞と数値は必ず一次資料で確かめる必要があります。}

\quizq{3}{大規模言語モデルが一般に苦手なこととして適切なものはどれでしょうか?}%
  {言語の翻訳}%
  {文章の要約}%
  {文体の変換}%
  {\correct{学習データ以降の最新情報を正確に答えること}}%
  {正解は (d)。LLM の知識には学習データの締め切り（知識カットオフ）があり、それ以降の出来事は知りません。Web検索や社内文書の検索と組み合わせる RAG（検索拡張生成）などで補います。(a)(b)(c) の翻訳・要約・文体変換は、LLM の得意分野です。}

\quizq{4}{「プロンプト」の説明として適切なものはどれでしょうか?}%
  {\correct{AIに与える指示や入力文}}%
  {AIの学習に使うデータセット}%
  {AIの処理速度の単位}%
  {AIが出力する文章}%
  {正解は (a)。プロンプトは AI に与える指示や質問の文です。(b) は学習データ、(d) は出力（応答）にあたります。本講義では、役割・目的と背景・条件・出力の形式という4要素を意識してプロンプトを書きます。}

\quizq{5}{「トークン」の説明として適切なものはどれでしょうか?}%
  {セキュリティのための合言葉}%
  {画像の解像度の単位}%
  {\correct{モデルが扱う文章の最小単位(語や部分語)}}%
  {AIの利用料金のこと}%
  {正解は (c)。トークンは LLM が文章を処理する最小単位で、単語や単語の一部に相当します。日本語では1文字が1〜2トークンになることが多いです。文脈窓（一度に扱える長さ）や API の料金はトークン数で決まりますが、(d) のようにトークンそのものが料金なのではありません。}

